\section{Conclusion}
\label{sec:conclusion}

This paper has provided a unified framework for the six partisan fairness metrics
that practitioners must navigate in post-\textit{Rucho} redistricting litigation.
The central contributions are three.

First, the taxonomy of descriptive vs.\ normative metrics clarifies what each
measure captures and what it cannot. Descriptive metrics (EG, MM, Bias, Declination)
characterize the partisan properties of a given map; they do not specify what the
properties should be. Normative metrics (proportionality, responsiveness) impose a
target; they require a prior normative commitment that courts must make explicitly.
Confusing these categories---treating a descriptive metric as if it contained its
own standard, or treating a normative metric as a value-free measurement---produces
avoidable expert testimony errors.

Second, the decision table (Table~\ref{tab:decision-table}) maps state constitutional
legal theory to metric selection, reducing the risk of metric-shopping criticism.
Expert witnesses should lead with the metric most directly relevant to the applicable
state constitutional standard, report all six for convergence, and explain the
relationships among them to preempt opposing arguments that any single number
was cherry-picked.

Third, the \textsc{bisect} reference distribution provides state courts with
a conceptually clear, evidentiary tool for measuring the magnitude of partisan
manipulation. Because \textsc{bisect} maps cannot be instructed to gerrymander---they lack the
partisan data inputs that gerrymandering requires---their partisan properties
reflect geographic sorting alone. That is a process safeguard, not a guarantee
of partisan neutrality. The gap between \textsc{bisect} values and enacted map
values on any metric measures the portion of partisan skew that is attributable
to deliberate mapmaker choices rather than residential patterns.
For North Carolina, this gap is 3.5$\times$ to 7$\times$ larger than what
geography alone produces, varying by metric; efficiency gap and mean-median
ratios are near 3.5$\times$, while partisan bias ratios reach 7$\times$ in
enacted maps with extreme manipulation.

The L.1--L.6 companion papers provide dedicated treatments of each metric,
including the full mathematical development, empirical results across NC, WI,
TX, and FL, legal history, and known limitations. Courts, commissioners, and
expert witnesses should use this overview paper to orient the framework and
consult the dedicated papers for the depth of analysis required in contested
proceedings.

\bigskip

\noindent\textit{The author thanks the redistricting research team for data
access and algorithmic development. All errors are the author's own.}
